
\section{Discussion}
\label{sec:discussion}

Our study answers the three questions posed in the introduction and turns them
into a design principle for agentic continual learning.

\paragraph{How should experience be consolidated?}
Experience should be transformed into shared, revisable procedural memory
rather than accumulated as independent traces or abstractions. By organizing
interlinked pages, \method{} can merge and correct knowledge across sessions
while preserving useful relationships at different granularities. Free Form
\method{} improves over No Memory by 10.9--14.5 pp across all six
model--benchmark settings with 92--255 KB stores---two orders of magnitude
smaller than raw-trace retrieval. The value of memory lies in selective
consolidation, not volume.

\paragraph{Which feedback, ingestion, and refinement choices matter?}
Failure traces are useful primarily when feedback makes them interpretable.
Execution-grounded outcomes localize failures, diagnosis explains the required
correction, and iterative refinement turns newly exposed gaps into durable
knowledge. One diagnostically supervised round outperforms four Self-Judge
rounds by 7.0--10.9 pp: feedback quality matters more than rollout volume.

\paragraph{Where should human effort be spent?}
General system design remains necessary: memory needs persistence, retrieval,
update operations, navigation, and quality checks. Domain-specific
organization, however, should be left to the agent rather than prescribed
before the knowledge gap is observed. Human effort is better spent reactively
diagnosing failures. Additional constraints should address demonstrated memory
failures, not anticipated ones.

\method{} therefore contributes both a training-free mechanism for acquiring
proprietary procedural knowledge and a practical principle: let the agent
organize memory, and spend human effort on diagnosis.

\paragraph{Limitations.}
Our controlled study uses two benchmarks, three models, and single-run
configurations; variance across seeds remains unquantified
(Appendix~\ref{app:limitations}). More broadly, enterprise procedural
automation remains less studied than code, mathematics, or personal assistants,
and no available open benchmark combines stateful environments, proprietary
procedures, and execution-grounded feedback at scale. Longer-term memory
stability, broader enterprise environments, replication, and scalable automated
diagnosis remain important directions for future work.
